\section{The absolute localization axiom and its limitations}
\label{sec:G}
\providecommand{\Wt}{\widetilde W}
\providecommand{\Wsharp}{W^{\sharp}}
\providecommand{\Hot}{\operatorname{Hot}}
\providecommand{\HOT}{\operatorname{HOT}}
\providecommand{\tors}{\operatorname{tors}}
\providecommand{\Gpd}{\operatorname{Gpd}}

Throughout this section ``Loc'' means the numbering of Chapter XVI unless
we explicitly write ``XV''.  In particular, a class satisfying only
Loc~(1)--(3) is called a \emph{weak localizer}; a \emph{fundamental
localizer} in the modern sense additionally satisfies relative Theorem~A,
Loc~(3 bis).  This distinction is essential here.
Our manuscript references are to the attached transcriptions
\cite{GDERXV,GDERXVI}; we check the supplied reviews
\cite{GR20,GREVLOC} where they enter the argument.

\begin{claim}{G.1}{Proved}{Axiom ledger and the test class of REVLOC F.1}
\label{G:Wt-axioms}
For a functor $v:A\to B$ between small categories, let $v\in\Wt$ mean
that the following four maps are isomorphisms, with an equivalence of
groupoids in the first line:
\begin{align*}
 &\Pi_1(BA)\longrightarrow\Pi_1(BB),\\
 &H_q(BA;\Z)\longrightarrow H_q(BB;\Z)\qquad(q\geq3),\\
 &H_2(BA;\Q)\longrightarrow H_2(BB;\Q),\\
 &\tors H_2(BA;\Z)\longrightarrow\tors H_2(BB;\Z).
\end{align*}
Then $\Wt$ is strongly saturated and satisfies the literal XVI
Loc~(1), (2), (3), (4), (5), (8), and (9).  Its aspherical categories
are exactly the nonempty weakly contractible categories.  REVLOC F.1's
Loc~(2), (3), (8), (9) already uses the XVI numbering; it omits, rather
than replaces, Loc~(1).
\end{claim}
\begin{proof}
The relevant numbering is as follows.  XVI Loc~(1) consists of inclusion
of isomorphisms, two out of three, and the rule that $gf,fg\in W$ imply
$f,g\in W$.  Loc~(2) concerns terminal objects, Loc~(3) is absolute
Theorem~A, Loc~(4) is arbitrary disjoint-union detection, Loc~(5) is
preservation of components, Loc~(8) is filtered-colimit closure, and
Loc~(9) is the small generating bound.  XV's corresponding numbers are
respectively (1), (2), (3), (5 bis), (8), (7), and (7 bis), with XV
(7 bis) combining the last two requirements.  XVI Loc~($6'$ a) is the
forward implication in the sieve-cover square, namely XV Loc~(4)'s
forward half.  None of these identifications uses a manuscript claim
about their consequences.

The four displayed tests are functorial.  For the first, we may pass
from groupoids to their category localized at equivalences, where
equivalences are exactly the isomorphisms.  Each test therefore satisfies
all of Loc~(1), including the rule about $gf$ and $fg$.  Each also
factors through $\Cat[\Wt^{-1}]$.  Consequently a functor inverted by
this localization passes every test and is already in $\Wt$.  This
proves strong saturation, and in particular retract closure.  All weak
homotopy equivalences belong to $\Wt$.

If $A\to e$ is in $\Wt$, then $BA$ is connected and simply connected.
Put $T=H_2(BA;\Z)$.  We have $T\otimes\Q=0$ and $\tors T=0$.
The first condition makes every element torsion, and the second kills
every element.  Hence $T=0$; homology in larger degrees vanishes by
definition.  A simply connected acyclic CW complex is weakly
contractible: a least nonzero homotopy group, if present, contradicts
Hurewicz.  Conversely a weakly contractible category passes the tests.
This checks the aspherical-object assertion of REVLOC F.1 directly.
It proves Loc~(2), and Quillen's Theorem~A
\cite[\S1, Theorem A]{GQuillen1973} proves Loc~(3), because its
hypotheses now say exactly that the comma categories are weakly
contractible.  The nerve of the empty category fails the first test
against $e$.

The fundamental-groupoid test gives a bijection of components.  Homology
and its torsion subgroup split as direct sums over components, while
fundamental groupoids split as disjoint unions.  Thus a disjoint union
of functors is in $\Wt$ if and only if every summand is.  This proves
Loc~(4) and (5), including the empty family.

For completeness, the two size assertions in REVLOC F.1 can be checked
without invoking an accessibility slogan.  A simplex in a nerve, a
path or relation between paths, an integral cycle or boundary, and a
rational denominator are all finite data.  They therefore commute with
filtered colimits of categories.  Torsion also commutes with filtered
colimits of abelian groups: an equation $nx=0$ in the colimit holds
at a later stage, and equality of torsion representatives does too.
These observations prove Loc~(8).

Given a member $v:A\to B$ of $\Wt$ and a countable subarrow, enlarge
its source and target by the witnesses for all countably many requests
for paths, lifts of paths, relations, homology lifts and boundaries,
rational lifts, and torsion lifts and boundaries.  In the last case we
also add a finite boundary witnessing torsion of the chosen source
class.  Each request has a finite witness in the given global arrow.
Close under composition and repeat countably many times.  The union is
a countable subarrow in $\Wt$, because every finite request occurs and
is answered.  The same construction beginning with a countable union
proves that these subarrows form a filtered cofinal family.  A set of
isomorphism representatives of countable arrows therefore generates
$\Wt$ by filtered colimits.  This is precisely the bounded-subcategory
form of XVI Loc~(9) used in REVLOC F.1.  We have checked all the stated
conditions, rather than assuming that the register's four numbers
include saturation.
\end{proof}

\begin{claim}{G.2}{Proved}{The finite failure of relative Theorem A}
\label{G:finite-relative}
The class $\Wt$ fails Loc~(3 bis) and Loc~($6'$ a).  The relative
Theorem~A witness in REVLOC F.1 has exactly $77$ source objects and
$53$ target objects, and its homological obstruction is integral
$2$-torsion.
\end{claim}
\begin{proof}
Let $C_m$ be the polygon with vertices $0,\ldots,m-1$ and edges
$\{i,i+1\}$, with indices modulo $m$.  Let
$L_m=C_m*\{a_m,b_m\}$, with no edge between $a_m,b_m$.  Thus $|L_m|$
is a $2$-sphere.  We use all nonempty faces, ordered by inclusion, to
form the finite posets
\[
 P=\operatorname{Face}(L_8),\qquad
 Q=\operatorname{Face}(L_4).
\]
They have $50$ and $26$ objects.  The simplicial map
$h:L_8\to L_4$, $h(i)=i\bmod4$, $h(a_8)=a_4$, $h(b_8)=b_4$,
induces a functor $f:P\to Q$ by taking the image of a face.  The
integral fundamental cycles
\[
 z_m=\sum_{i=0}^{m-1}
 \bigl([a_m,i,i+1]-[b_m,i,i+1]\bigr)
\]
satisfy $h_*z_8=2z_4$.  Hence $Bf$ has degree two.  It lies in
$\Wt$: both spaces are simply connected, their homology above degree
two vanishes, their $H_2$ is torsion free, and multiplication by two is
invertible over $\Q$.

Let $J$ be the poset with arrows $0\to1$, $0\to2$ and no other
nonidentity arrows.  Consider covariant diagrams in $\Cat$
\[
 D=(Q\xleftarrow{f}P\to e),\qquad
 E=(Q\xleftarrow{\id}Q\to e).
\]
The natural transformation $D\to E$ has components $f,\id_Q,\id_e$.
Its opfibrational Grothendieck construction is $v:S\to T$ over $J$.
Explicitly, $S$ has the $P$ layer, the $Q$ layer, and an object $t$;
the original orders remain, $p\leq q$ means $f(p)\leq q$ in $Q$,
and $p\leq t$ for every $p\in P$.  There are no comparisons between
$Q$ and $t$.  Replace $P,f$ by $Q,\id_Q$ to obtain $T$.  All arrows
and all compositions have now been specified.

For each $j\in J$, evaluation
\[
 S/j\longrightarrow D(j),\qquad(i,x,\alpha:i\to j)
 \longmapsto D(\alpha)x
\]
is left adjoint to the inclusion $x\mapsto(j,x,\id_j)$, and similarly
for $E$.  These adjunctions commute with $v/j$ and the displayed
component maps.  Thus each $v/j$ is in $\Wt$.

The full sieves $P\cup Q$ and $P\cup\{t\}$ cover $S$, and intersect
in $P$.  Every simplex in $NS$ lies in one of these pieces.  Consequently
their realizations give an actual union of CW subcomplexes and hence a
homotopy-pushout square.  The first piece retracts by an adjunction onto
$Q$; the second is a cone on $P$.  Thus $BS$ is the homotopy cofiber
of $Bf$, and $BT$ is the cofiber of an identity.  The former is simply
connected and has reduced integral homology $\Z/2$ in degree two only;
the latter is contractible.  In particular $v\notin\Wt$.

In the same sieve-cover square, $P\to P\cup Q$ is in $\Wt$ because
its composite with the adjoint retraction is $f$.  But
$P\cup\{t\}\to S$ is not in $\Wt$, by the just computed torsion
group.  This refutes Loc~($6'$ a).

The independent program \texttt{code/G\_finite\_witnesses.py} constructs
the orders and functors, verifies the three evaluation adjunctions and
the sieve cover, and checks the integral chain equality above.  It
finds $521$ arrows in $S$ and $321$ in $T$, including identities.
Its cofiber certificate is the integer differential $[2]$ in degrees
three and two, compared with $[1]$ for the target.  Thus the obstruction
is not inferred from rational ranks.
\end{proof}

\begin{claim}{G.3}{Refuted}{Products under the weak axioms}
\label{G:products-fail}
Neither of the assertions
\[
 f\in W\Longrightarrow f\times B\in W,
 \qquad \Cat\longrightarrow\Hot_W\text{ preserves finite products}
\]
follows from XVI Loc~(1)--(3).  Both fail for the strongly saturated
class $W=\Wt$, already with finite posets.
\end{claim}
\begin{proof}
Use $f:P\to Q$ from \ref{G:finite-relative}, and let
$B=\operatorname{Face}(C_4)$, an eight-object circle poset.  The functor
$f\times\id_B:P\times B\to Q\times B$ has $400$ source objects
and $208$ target objects.  Integral K\"unneth, or the tensor product of
the displayed cycle complexes, gives
\[
 H_3(B(P\times B);\Z)\cong\Z,
 \qquad H_3(B(f\times\id_B);\Z)=[2].
\]
Thus it is not in $\Wt$.  The program constructs the complete product
orders ($3104$ and $1568$ arrows) and the product functor.  On cellular
product chains it checks exactly
$(h\times\id)_*(z_8\times c_4)=2(z_4\times c_4)$.

If the localization functor preserved binary products, it would carry
$f\times\id_B$ to the product of an isomorphism and an identity,
therefore to an isomorphism.  Strong saturation in
\ref{G:Wt-axioms} would then put this functor in $\Wt$, a contradiction.
This refutes preservation of the specified products, not merely a
particular proposed proof of preservation.
\end{proof}

\subsection{The nontriviality propositions require care with saturation}

\begin{claim}{G.4}{Proved}{A weak localizer with component merging}
\label{G:Wsharp}
There is a class $\Wsharp$ satisfying XVI Loc~(1)--(4), with exactly
the weakly contractible aspherical categories, which contains a
functor merging two components but for which $\Hot_{\Wsharp}$ has
infinitely many isomorphism classes.  In particular strong saturation
cannot be silently added to the weak axioms.
\end{claim}
\begin{proof}
For a small category $X$, write $X=X_0\amalg X_+$, where $X_0$ is the
union of its connected components $C$ with $H_1(BC;\Z)=0$, and $X_+$
is the union of all other components.  For a functor $u:X\to Y$ put
$u\in\Wsharp$ precisely when the following requirements hold:
\begin{enumerate}
\item $u$ preserves this division in both directions: a component of
$X$ lies in $X_0$ if and only if its image lies in $Y_0$;
\item $u_0:X_0\to Y_0$ is a weak homotopy equivalence;
\item for each connected component $D\subseteq Y_+$ the specified
homomorphism
\[
 \bigoplus_{\substack{C\in\pi_0(X_+)\\u(C)\subseteq D}}
 H_1(BC;\Z)\longrightarrow H_1(BD;\Z)
\]
is an isomorphism.
\end{enumerate}
Empty sums mean the zero group.  Since each group in the target of
the last line is nonzero, every $Y_+$ component has a preimage.
The second requirement supplies a bijection on the $0$ components.
Thus every member of $\Wsharp$ is surjective on components, and
induces an isomorphism on global integral $H_1$.

Here is a full verification of two out of three.  For composable
$X\xrightarrow{u}Y\xrightarrow{v}Z$ preserving the division,
the last condition can equivalently be tested on the direct sum over
all positive components: the map is block diagonal, and a block
diagonal map is an isomorphism exactly when each block is.  If $u,v$
are in $\Wsharp$, all three conditions compose.  If $v,vu$ are in
$\Wsharp$, preservation of the division by $v$ and $vu$ forces it
for $u$.  Two out of three on global $H_1$ gives the positive condition
for $u$, and two out of three on the zero parts gives the second
condition.  If $u,vu$ are in $\Wsharp$, surjectivity of $\pi_0(u)$
forces $v$ to preserve the division.  The same two arguments then
give the remaining conditions for $v$.

We check also the literal Loc~(1)(c), which is stronger than just
checking idempotent splittings.  Suppose $vu$ and $uv$ are in
$\Wsharp$.  The two global $H_1$ composites are isomorphisms, so
$H_1(u)$ and $H_1(v)$ are isomorphisms.  A positive component cannot
be sent by $u$ to a zero component: its nonzero $H_1$ summand would
then be killed by $vu$, whereas the last condition makes its
restriction injective.  If a zero component were sent by $u$ to a
positive one, preservation of the division by $vu$ would send this
positive component back to a zero component under $v$; the composite
$uv$ would then kill a nonzero $H_1$ summand.  This is impossible too.
The same argument applies to $v$.  On the zero parts, both composites
are weak homotopy equivalences, so $u_0,v_0$ are weak homotopy
equivalences.  The global $H_1$ isomorphisms give the positive blocks.
This proves (c).  Isomorphisms plainly belong to the class.

All weak homotopy equivalences belong to $\Wsharp$.  The map $X\to e$
belongs to it if and only if $X=X_0$ and $X\to e$ is a weak homotopy
equivalence.  Thus the aspherical objects are precisely the weakly
contractible ones.  Terminal objects and absolute Theorem~A therefore
give Loc~(2) and (3).  All three defining tests operate separately on
target components, so they detect arbitrary disjoint unions in both
directions; this is Loc~(4).

For the promised finite functor, take the face poset $A$ of two
disjoint triangular circles and the face poset $B$ of their wedge at
one vertex.  Map each circle to its corresponding wedge summand.
The posets have respectively $12$ and $11$ objects.  Their positive
component homology is $(\Z,\Z)$ and $\Z^2$, and the displayed
homomorphism is the identity matrix.  Thus $a:A\to B$ lies in
$\Wsharp$ although $\pi_0(a):\{0,1\}\to\{*\}$ is not injective.
The program specifies the six edges of each graph, all face-poset
arrows and the vertex identification, and checks the integral matrix.

Since $H_1(-;\Z):\Cat\to\Ab$ inverts $\Wsharp$, it factors through
$\Hot_{\Wsharp}$.  Finite wedges of $r$ circles have homology $\Z^r$
and hence give infinitely many isomorphism classes in that localization.
Also the identity and a constant endofunctor of a circle poset have
different images there.  Finally, $\{0,1\}\to e$ is an arrow-category
retract of $a$: select the two vertices identified by $a$, and use the
component map $A\to\{0,1\}$ and the unique map $B\to e$.  The
retraction squares commute strictly.  It is therefore inverted after
localization, although it is not in $\Wsharp$.  This exhibits the
failure of strong saturation and of arrow-retract closure explicitly.
\end{proof}

\begin{claim}{G.5}{Refuted}{XV 9.11, 9.12, and 9.13 under the literal weak axioms}
\label{G:nontriviality-false}
Each of XV Proposition~9.11, Proposition~9.12, and Corollary~9.13,
with only XVI Loc~(1)--(3) as hypotheses, is false.  For the particular
test class $\Wt$, however, all the conditions of 9.11 and 9.12 hold,
and 9.13 holds vacuously.
\end{claim}
\begin{proof}
Use $\Wsharp$ from \ref{G:Wsharp}.  In 9.11, condition (b) is true:
arbitrary coproduct detection holds, and the homotopy category has
infinitely many isomorphism classes.  Conditions (a) and (c) are false,
because $a:A\to B$ merges components.  This is a counterexample to
the equivalence asserted in that proposition, including its
decomposition hypothesis.

In 9.12, (a) is false.  So is (b): let $x_0,x_1:e\to A$ select the
two vertices identified by $a$.  We have $ax_0=ax_1$ strictly.  Since
$a$ is inverted, the two maps $x_0,x_1$ are equal in the localization,
although their component maps differ.  Condition (c) is true, because
the identity and the constant endofunctor of a circle induce different
maps on $H_1$.  Conditions (c'), (c''), and (d) are true by the
infinitely many homology ranks in \ref{G:Wsharp}.  The same identity
and constant endofunctor refute 9.13, whose premise is true for
$\Wsharp$.

For $\Wt$, the first defining test preserves components and all four
tests detect coproducts.  Hence 9.11(a),(c) hold.  The functor $\pi_0$
factors through the localization, so finite discrete categories of
different cardinalities remain nonisomorphic; 9.11(b) and 9.12(d)
follow.  It also gives 9.12(b) and distinguishes the two maps
$e\rightrightarrows\{0,1\}$, proving (c), (c'), and (c'').
Condition 9.12(a) is already known.  Since the failure of the
component axiom never occurs for $\Wt$, the implication 9.13 is
vacuous for this class.  These are explicit verifications for the
required test class, not counterexamples wrongly assigned to it.
\end{proof}

\begin{claim}{G.6}{Proved}{What survives and the direction of the edition's note}
\label{G:ccprime}
For every $W$ satisfying XVI Loc~(1)--(3), retain the six conditions
(a),(b),(c),(c'),(c''),(d) of XV~9.12 verbatim.  One always has
\[
 (a)\Longrightarrow(b)\Longrightarrow(c)
 \Longleftrightarrow(c')\Longleftrightarrow(c''),
 \qquad (a)\Longrightarrow(d)\Longrightarrow(c').
\]
In particular (c)$\Rightarrow$(c') is immediate, and
(c')$\Rightarrow$(c) is proved by the quotient-category argument at
XV manuscript pages 133 ter--133 quater.  The edition's note at
133 bis reverses these directions when read for the affirmative
conditions printed in 9.12.  The corrected unconditional version of
9.13 is the exact equivalence
\[
 \bigl(\text{every }\Hom_{\Cat}(X,Y)\to\Hom_{\Hot_W}(X,Y)
          \text{ is constant}\bigr)
 \Longleftrightarrow \Hot_W^*\simeq e.
\]
For XV~9.11, (a)$\Longleftrightarrow$(c)$\Longrightarrow$(b) always
holds; its converse (b)$\Rightarrow$(a) fails as in \ref{G:nontriviality-false}.
\end{claim}
\begin{proof}
If $W$ preserves components, $\pi_0$ factors through the localization.
This proves (a)$\Rightarrow$(b), and also (a)$\Rightarrow$(d) by
using discrete sets of all finite cardinalities.  Taking the two
point inclusions into a discrete two-point category proves
(b)$\Rightarrow$(c).  A nonconstant comparison on a hom set must have
nonempty source and target, because an empty source has only one
actual functor to any category and an empty target has at most one.
It therefore witnesses (c') immediately.

For the converse, assume every comparison on actual hom sets is
constant.  If $X$ is nonempty, choose $x:e\to X$ and write $p:X\to e$.
The two actual endofunctors $xp$ and $\id_X$ become equal; since
$px=\id_e$, $X$ becomes isomorphic to $e$.  Moreover every actual
functor between nonempty categories becomes the same canonical
morphism.  Every localized morphism is a finite composite of actual
functors and formal inverses of members of $W$; hence the entire full
subcategory on the nonempty categories has one morphism between
any two objects.  This proves $\Hot_W^*\simeq e$.

Here is the quotient argument with its ambiguity removed.  Let
$\mathcal H$ have the same objects as $\Cat$, one arrow between
nonempty objects, one arrow from $\emptyset$ to every object, and no
arrow from a nonempty object to $\emptyset$.  It is equivalent to
$[1]$.  The preceding hypothesis factors the localization functor
through $\mathcal H$.  If no member of $W$ goes from the empty category
to a nonempty one, the universal property of localization identifies
$\Hot_W$ with this quotient up to equivalence.  If such a member
exists, localizing $\mathcal H$ at it gives the category equivalent
to $e$.  Thus $\Hot_W^*\simeq e$ implies the negation of (c'').
The converse needs a small observation under the weak axioms.  If
$\emptyset\to X$ belongs to $W$ with $X$ nonempty, every endofunctor
of $X$ belongs to $W$ by two out of three.  Choose a point of $X$;
the constant endofunctor and Loc~(1)(c) imply $X\to e\in W$, hence
$\emptyset\to e\in W$.  Absolute Theorem~A applied to
$\emptyset\to Y$ then makes every category aspherical and every
functor a member of $W$.  Otherwise no zigzag in the localization can
go from a nonempty category to the empty one, since its first such
crossing would invert a member of $W$ with empty source.  Therefore
if $\Hot_W\simeq[1]$, all nonempty objects must lie in its terminal
isomorphism class.  This proves (c')$\Longleftrightarrow$(c'').
The implication (d)$\Rightarrow$(c') is now immediate.

This is exactly the missing contrapositive step after the statement
of 9.13.  Grothendieck's phrase that
``$(b)\Longrightarrow(c)\Longrightarrow(c')$ sont elles aussi triviales''
has the correct direction for the stated conditions.  We disagree on
that direction with the edition's note, Chapter XV, ms.~133 bis
(K\"unzer, Malgoire, Maltsiniotis); we agree with its location of the
nontrivial argument at ms.~133 ter--quater.  We also correct the
transcription's target of $h=(f_0,f_1)$ at ms.~133 quinquies:
it is $X\amalg X\to Y$, not $X\amalg X\to X$.

Finally, 9.11(a) says exactly component preservation together with
arbitrary coproduct detection.  Decompose any functor by target
components to obtain (c).  Conversely, (c) directly gives component
preservation; applying it to a disjoint union of functors gives
arbitrary coproduct detection.  Either condition implies (b), since
component preservation gives infinitely many isomorphism classes.
\end{proof}

\begin{claim}{G.7}{Proved}{Retract closure restores all three nontriviality assertions}
\label{G:retract-restoration}
If $W$ satisfies XVI Loc~(1)--(3) and is closed under retracts in the
arrow category, then XV~9.11, 9.12, and 9.13 hold in full.  In fact,
either $W\subseteq W_0$, or $W$ contains every functor between
nonempty categories; in the latter case $\Hot_W$ is equivalent to
$[1]$ or $e$.  Thus Loc~(1 bis), strong saturation, or Loc~(8) each
suffices for these three assertions, without product stability.
\end{claim}
\begin{proof}
First suppose $u:X\to Y$ in $W$ is not injective on components.
Choose $x_0,x_1$ in different components whose images are joined by
a finite zigzag in $Y$.  We can add a finite whisker to the component
of $x_1$, one fresh object and one arrow at a time, so that $u$ extends
to $u':X'\to Y$ and the end of the whisker maps exactly to $u(x_0)$.
At each step the old category is a deformation retract: fold the
new object to the attachment object; the added arrow gives a natural
transformation between the identity and this fold-inclusion composite.
Thus $X\to X'$ is in $W_\infty\subseteq W$, and $u'\in W$ by two
out of three.  The inclusion $W_\infty\subseteq W$ is Cisinski's
Theorem~6.1.18 \cite{GCisinski2006}, whose weak-saturation axiom is implied by the literal
Loc~(1); no relative localization axiom occurs in that theorem.

Select the two objects of $X'$ with the same image, and map $X'$ to
the discrete two-point category by its components, separating those
objects.  Together with $e\to Y$ selecting their common image and
$Y\to e$, these give a strict arrow-category retraction of $u'$
onto $p:\{0,1\}\to e$.  Hence $p\in W$.

If instead $u$ is not surjective on components and $X\ne\emptyset$,
choose a point of $X$ and a point of a missing target component.
Map $Y$ to $\{0,1\}$ by assigning all image components to $0$ and
the chosen missing component to $1$.  Then $e\to\{0,1\}$ is an
arrow-category retract of $u$.  It lies in $W$, and two out of three
puts $p$ in $W$.  If $X$ is empty, the observation in
\ref{G:ccprime} already makes $W$ the class of all functors.

Assume now $p\in W$.  For every category $A$ the projection
$A\times\{0,1\}\to A$ belongs to $W$ by absolute Theorem~A:
its slice at $a$ is $(A/a)\times\{0,1\}$, whose projection to
$\{0,1\}$ is a weak homotopy equivalence, since $A/a$ has a terminal
object.  Both endpoint sections are in $W$.  If $A$ is nonempty,
choose $a:e\to A$ and let $r:A\to e$.  The functor
$A\times\{0,1\}\to A$ restricting to $\id_A$ and $ar$ shows by
two out of three that $ar\in W$.  Since $ra=\id_e$, Loc~(1)(c)
gives $a,r\in W$.  Every nonempty category is therefore aspherical,
and every functor between nonempty categories is in $W$.

This dichotomy, with \ref{G:ccprime}, proves every equivalence in
9.12 and all of 9.13; it also restores the missing implication of
9.11.  Lastly, filtered-colimit closure implies arrow-retract closure.
Indeed if an arrow $a$ is a retract of an arrow $b\in W$, the
corresponding idempotent endomorphism of $b$ defines a constant-object
sequence $b\to b\to\cdots$ in $\Cat^{[1]}$ whose colimit is the
split image $a$.  Loc~(8) puts this colimit in $W$.  This explains why
the three propositions hold for $\Wt$ even though products fail there.
\end{proof}

\subsection{A left adjoint need not exist}

\begin{claim}{G.8}{Refuted}{Failure of the bare existence of a homotopy left Kan extension}
\label{G:no-left-adjoint}
Let $J=(0\to1,0\to2)$, let $I=J^{\op}$, and let $u:I\to e$.
For $W=\Wt$ the functor
\[
 u^*:\Hot_{\Wt}\longrightarrow
 \Cat^{I^{\op}}[(\Wt)^{I^{\op}}{}^{-1}]
       =\HOT_{\Wt}(I)
\]
of constant diagrams has no left adjoint.  Thus the bare existence
claim in XV~10.1, $1^\circ$, fails under XVI Loc~(1)--(3), even when
strong saturation and Loc~(8),(9) are added.
\end{claim}
\begin{proof}
We give the obstruction in the indicated localized categories, so
failure of the usual Grothendieck-construction formula is not being
substituted for nonexistence of an adjoint.

For a finite poset $R$, let $D(R):J\to\Cat$ be the span
$e\leftarrow R\to e$.  Let $K(R)$ be the poset obtained from $R$
by adding two incomparable objects, each above every object of $R$.
Then $BK(R)$ is the unreduced suspension of $BR$.  The diagram
$\overline D(R)$ replaces each endpoint $e$ by a separate terminal
cone $R^{\triangleright}$, with its evident inclusion from $R$.
There is a natural objectwise weak homotopy equivalence
$r_R:\overline D(R)\to D(R)$ and a strict natural cocone
$c_R:\overline D(R)\to u^*K(R)$.  Write
\[
 \alpha_R=c_Rr_R^{-1}:D(R)\longrightarrow u^*K(R)
\]
in $\HOT_{\Wt}(I)$.  These maps are natural in $R$.

Use $P,Q,f$ from \ref{G:finite-relative}.  In addition define the
simplicial degree-one map $k:L_8\to L_4$ on polygon vertices by
\[
 (0,1,2,3,4,5,6,7)\longmapsto(0,1,2,3,0,0,0,0),
\]
and preserve the two suspension vertices.  The same letter denotes
the induced face-poset functor $k:P\to Q$.  The integer cycle
calculation gives $k_*z_8=z_4$, so $k$ is a weak homotopy equivalence
between $2$-spheres.  In the diagram localization define
\[
 a=D(f)D(k)^{-1}:D(Q)\longrightarrow D(Q).
\]
It is an automorphism, since $D(f)$ and $D(k)$ are objectwise in
$\Wt$.  In $\Hot_{\Wt}$ define
\[
 b=K(f)K(k)^{-1}:K(Q)\longrightarrow K(Q).
\]
Here $K(k)$ is a weak homotopy equivalence, whereas $H_3(b;\Z)$ is
multiplication by two on $H_3(BK(Q);\Z)=\Z$.  Naturality of the
finite cocones gives the exact equality
\[
 \alpha_Q a=u^*(b)\alpha_Q
       \quad\text{in }\HOT_{\Wt}(I).
\]

There is a useful invariant of this diagram localization:
\[
\begin{aligned}
 T(F)&=H_3\operatorname{Cone}(d_F),\\
 d_F&=(F_{01*},-F_{02*}):C_*(NF(0);\Q)\\
 &\hspace{45mm}\longrightarrow
 C_*(NF(1);\Q)\oplus C_*(NF(2);\Q),
\end{aligned}
\]
where $F:J\to\Cat$, $C_*$ denotes normalized simplicial chains,
and $\operatorname{Cone}$ is the homological mapping cone.
Objectwise members of $\Wt$ are rational homology equivalences in
every degree: this is explicit in degree two, follows from integral
homology in larger degrees, and follows from the fundamental-groupoid
condition in degrees zero and one.  The long exact sequence of a
mapping cone therefore shows that $T$ inverts them, so $T$ is well
defined on $\HOT_{\Wt}(I)$.  On a constant diagram the chain map is
$(\id,-\id):C_*(NZ;\Q)\to C_*(NZ;\Q)^{\oplus2}$; summing the
two copies gives a natural quasi-isomorphism of its cone with
$C_*(NZ;\Q)$.  Thus $T(u^*Z)=H_3(BZ;\Q)$ naturally in the
localized category.  Moreover the union of the two nerve cones on
$NQ$ is exactly $NK(Q)$.  Its chain sequence
\[
0\longrightarrow C_*(NQ;\Q)
\longrightarrow C_*(NQ^{\triangleright};\Q)^{\oplus2}
\longrightarrow C_*(NK(Q);\Q)\longrightarrow0
\]
is exact, with signs $(\id,-\id)$ on the first map.  Consequently
the cocone $c_Q$ induces a quasi-isomorphism from that mapping cone
to $C_*(NK(Q);\Q)$.  Since $r_Q$ is objectwise a weak equivalence,
$T(\alpha_Q)$ is an isomorphism $\Q\to\Q$.

Suppose a left adjoint $L$ to $u^*$ existed.  Put $Z=LD(Q)$ and let
$\eta:D(Q)\to u^*Z$ be its unit at $D(Q)$.  Let
$\beta:Z\to K(Q)$ correspond under the adjunction to $\alpha_Q$.
Then
\[
 \alpha_Q=u^*(\beta)\eta,
 \qquad \beta L(a)=b\beta.
\]
Integral $H_3$ factors through $\Hot_{\Wt}$ by definition of
$\Wt$.  Hence $H_3(L(a))$ is an automorphism of $H_3(BZ;\Z)$,
and the second equality gives
\[
 \operatorname{im}H_3(\beta)
       =2\operatorname{im}H_3(\beta)\subseteq\Z.
\]
The only subgroup of $\Z$ equal to twice itself is zero.  Therefore
$H_3(\beta;\Z)=0$, and also $H_3(\beta;\Q)=0$.  Applying $T$ to
the first equality now makes $T(\alpha_Q)=0$, contradicting the
computed isomorphism.  No left adjoint exists.

Every input category and actual functor in this obstruction is finite.
The program verifies both degrees, all face-poset maps and both
suspension maps; $K(P)$ and $K(Q)$ have $52$ and $28$ objects.
The formal inverses used above are legitimate inverses of the named
weak equivalences in the specified localizations.
\end{proof}

\subsection{The additional hypotheses which restore the familiar theory}

\begin{claim}{G.9}{Known}{Forward gluing and filtered colimits imply the relative axiom}
\label{G:gluing-recovery}
Let $W$ satisfy XVI Loc~(1)--(3).  If it also satisfies
Loc~($6'$ a) and Loc~(8), then it is a fundamental localizer in the
modern sense and satisfies Loc~(3 bis).  Loc~(9) is unnecessary.
This is REVLOC F.1.1, with its reduction to the opened results of
Cisinski, \emph{Ast\'erisque} 308 \cite{GCisinski2006}, checked below.
\end{claim}
\begin{proof}
Cisinski's Theorem~6.1.18 gives $W_\infty\subseteq W$ under exactly
the weak axioms.  Let $i_\Delta X=\Delta/X$ be the simplex category
of a simplicial set $X$, and define
\[
 \mathcal V=\{g:X\to Y\mid i_\Delta g\in W\}.
\]
It has two out of three and contains all weak homotopy equivalences,
in particular every simplicial map having the right lifting property
against all monomorphisms.

Suppose $X\hookrightarrow Y$ is a monomorphism in $\mathcal V$,
$X\to Z$ is arbitrary, and $T=Y\amalg_X Z$.  Put
$A=\Delta/X$, $B=\Delta/Y$, $C=\Delta/Z$, and form the double
collage $P$ of $B\leftarrow A\to C$.  Its full subcategories
$M_{A\to B}=A\cup B$ and $M_{A\to C}=A\cup C$ are sieves
covering $P$ and intersecting in $A$.  In each collage mixed arrows
$a\to b$ are the arrows from the image of $a$ to $b$ in the target;
there are no arrows in the reverse direction or between $B,C$.
Retraction of each collage onto its target is a left adjoint.
Thus $A\to M_{A\to B}$ lies in $W$, and forward gluing gives
$M_{A\to C}\to P\in W$.  The inclusion $C\to M_{A\to C}$ is
a weak homotopy equivalence; hence $C\to P\in W$.

The compatible functors from the three layers to $\Delta/T$ give
$P\to\Delta/T$.  It is a weak homotopy equivalence: on nerves the
two collages are subcomplexes with common subcomplex $NA$, so $BP$
is the homotopy pushout of $BB\leftarrow BA\to BC$; the natural
simplex-category comparisons identify this with the homotopy pushout
of $|Y|\leftarrow|X|\to|Z|$.  Since $X\to Y$ is a monomorphism,
the simplicial pushout is a homotopy pushout.  It follows that
$\Delta/Z\to\Delta/T$ lies in $W$.  We have proved pushout closure
of monomorphisms in $\mathcal V$ without relative Theorem~A.

The functor $i_\Delta$ commutes with filtered colimits: its objects,
arrows, and the one equation defining an arrow are finite data.
Loc~(8) therefore makes $\mathcal V$ closed under filtered colimits.
Composition and the limit-ordinal step prove closure of its
monomorphisms under arbitrary transfinite compositions.  These are
exactly the three axioms in Cisinski's Definition~1.4.1 for a
$\Delta$-localizer.  Proposition~3.4.25 says that every such
localizer is regular.  Its representables are already weakly
contractible, so it is a test localizer in Definition~4.2.11.
Theorem~4.2.15 associates a fundamental localizer $W'$ on $\Cat$
such that $\mathcal V=i_\Delta^{-1}W'$.  We can identify it without
confusing the right adjoint to the simplex-category construction with
the ordinary nerve.  The natural last-vertex functor
$\Delta/NA\to A$ is a weak homotopy equivalence, and both $W$ and
$W'$ contain $W_\infty$.  Its square for any $v:A\to B$ and two out
of three yield
\[
 v\in W\ \Longleftrightarrow\ i_\Delta Nv\in W
 \ \Longleftrightarrow\ Nv\in\mathcal V.
\]
The same square gives
$v\in W'\Longleftrightarrow i_\Delta Nv\in W'
\Longleftrightarrow Nv\in\mathcal V$.
Hence $W=W'=N^{-1}\mathcal V$ and satisfies the relative axiom;
this also verifies directly the inverse-image identification used in
REVLOC F.1.1 and in Remark~4.2.16.
We have opened the precise Cisinski statements, including the
weak-saturation definition; no accessibility hypothesis has been
inserted into them.
\end{proof}

\begin{claim}{G.10}{Proved}{Restored products and left Kan extensions}
\label{G:restored-theory}
Under XVI Loc~(1), (2), (3 bis), the class $W$ is stable under finite
products, $\Cat\to\Hot_W$ preserves them, and every restriction
$u^*:\HOT_W(J)\to\HOT_W(I)$ has a left adjoint for a functor
$u:I\to J$.  In particular the failures of \ref{G:products-fail}
and \ref{G:no-left-adjoint} are repaired, for a class already
satisfying Loc~(8) such as $\Wt$, by adding Loc~($6'$ a).  Loc~(8)
alone does not repair either failure.
\end{claim}
\begin{proof}
Let $f:A\to A'$ lie in $W$.  Regard $f\times\id_B$ as a functor
over $B$.  Its slice at $b$ is $f\times\id_{B/b}$.  The two
projections from $A\times(B/b)$ and $A'\times(B/b)$ to $A,A'$
are homotopy equivalences because $B/b$ has a terminal object.
Thus the sliced functor lies in $W$ by two out of three.  Relative
Theorem~A puts $f\times\id_B$ in $W$.  Binary products of arbitrary
members follow by composition, and the empty product is an identity.
The purely categorical finite-product localization result in
\ref{G:finite-products-general} below then gives preservation.

We spell out the Kan-extension comparison, including its variance.
For covariant diagrams on a small category $K$, let $W_K$ denote
objectwise membership in $W$, and let $W'_K$ be the class of functors
over $K$ whose slices at every object lie in $W$.  There is an
adjunction between $\Cat/K$ and $\Cat^K$: one functor sends
$p:A\to K$ to $k\mapsto A/k$, and the other sends a diagram to its
opfibrational Grothendieck construction.  The counit at a diagram
$F$ is evaluation
$(\int_KF)/k\to F(k)$, left adjoint to inclusion of the $k$-fiber.
It is objectwise a weak homotopy equivalence.  The unit has the
corresponding local property.  The adjunction therefore gives
equivalences
\[
 (\Cat/K)[(W'_K)^{-1}]\simeq\Cat^K[W_K^{-1}].
\]
Only the absolute axioms are used so far; this is the distinction
made in Maltsiniotis, \emph{La th\'eorie de l'homotopie de
Grothendieck}, Theorem~3.1.4 \cite{GMaltsiniotis2020}.

For $w:K\to L$, composition $\Cat/K\to\Cat/L$ respects these
local equivalences by the relative axiom.  Indeed at $\ell\in L$
the functor is over $K/\ell$, and slicing further at
$(k,wk\to\ell)$ recovers the original slice at $k$.
Transporting composition across the displayed equivalences gives
$w_!$.  For precision, the adjunction we descend is the following
cross-adjunction, not an unproved assertion that strict pullback
preserves local equivalences.  Its left functor sends $p:A\to K$
to the covariant diagram $\ell\mapsto A/\ell$ over $L$, with $A$
viewed over $L$ by $wp$; its right functor sends $F:L\to\Cat$ to
$\int_K Fw\to K$.  Its counit at $F$ is the natural evaluation
\[
 (\textstyle\int_KFw)/\ell\longrightarrow F(\ell),
 \qquad(k,x,\alpha:wk\to\ell)\longmapsto F(\alpha)x.
\]
Its unit at $p:A\to K$ sends $a$ to
$(p(a),a,\id_{wp(a)})$ in the Grothendieck construction of
$k\mapsto A/wk$.  These formulas on arrows use composition in $A,L$
and the action of $F$; substitution verifies both triangle identities.
The left functor is relative by the preceding relative-Theorem-A
argument, and the right one is relative by the already established
comparison between diagrams and local slice equivalences.  The
unit and counit therefore descend.  Transport through the equivalence
for $K$ makes this precisely $w_!\dashv w^*$.
This is the proof of Maltsiniotis's Lemma~3.1.5 and
Theorem~3.1.7; the source explicitly strengthens ``weak localizer''
to ``fundamental localizer'' immediately before Lemma~3.1.5.

To apply this to the prompt's contravariant diagrams, take
$K=I^{\op}$, $L=J^{\op}$, and $w=u^{\op}$.
For $F:I^{\op}\to\Cat$, a representative at $j\in J$ is the
opfibrational Grothendieck construction over
$(u^{\op}/j)\cong(j\backslash u)^{\op}$ of the restricted
covariant diagram.  This states the formula with its actual variance;
its equivalent fibred form is obtained by taking opposites of the
fiber categories as well.  The final assertion follows from
\ref{G:gluing-recovery} and the fact that $\Wt$ already satisfies
Loc~(8).  The counterexamples show directly that filtered-colimit
closure without forward gluing cannot suffice.
\end{proof}

\subsection{Localization and infinite products}

\begin{claim}{G.11}{Known}{The finite-product part of XVI Lemma 2.5}
\label{G:finite-products-general}
Let $\mathcal M$ be a category, let $W$ contain all identities,
and suppose that products indexed by a finite set $I$ exist in
$\mathcal M$ and that products of $I$-indexed families in $W$ lie
in $W$.  If $I$ is nonempty, it is enough to assume the corresponding
binary-product closure and identities.  The localization
$Q:\mathcal M\to\mathcal M[W^{-1}]$ preserves these products,
which therefore exist in the localization.  A terminal object is
always preserved, without a product-closure hypothesis.
\end{claim}
\begin{proof}
For finitely many relative categories, localization commutes with
their product.  A tuple of zigzags can be put into a common finite
alternating shape by adding identity arrows.  Given an equality of
two tuples, the finitely many finite derivations from the defining
localization relations can likewise be performed one coordinate at
a time.  Thus the canonical functor from the localization of the
product to the product of the localizations is full and faithful,
and is the identity on objects.

The adjunction
$\operatorname{diag}:\mathcal M\rightleftarrows\mathcal M^I:\prod_I$
has both functors relative under the hypotheses.  Localize its unit
and counit; their naturality and triangle identities remain valid.
Using the preceding identification of the localized product category
shows that the descended right adjoint is a product functor and that
the original product cone is preserved.  The same argument with the
one-object target category proves preservation of a terminal object.
Alternatively a zigzag ending at a terminal object reduces uniquely
to the terminal map using its defining equalities.

The finite statement is also formalized in Mathlib's
\texttt{CategoryTheory.\allowbreak Localization.\allowbreak FiniteProducts}, theorem
\texttt{preservesProductsOfShape}, whose hypotheses explicitly
include finiteness of the indexing type \cite{GMathlibFiniteProducts}.
We use that opened source
only as a literature record; the argument above supplies the proof.
\end{proof}

\begin{claim}{G.12}{Refuted}{The infinite-product assertion of XVI Lemma 2.5}
\label{G:infinite-product-counterexample}
There is a small category $\mathcal M$ with all countable products
and a class $W$ containing identities and stable under countable
products, such that the localization functor fails to preserve a
countable product.  Thus the infinite assertion of XVI ms.~7--16
is false in the stated generality.
\end{claim}
\begin{proof}
Let $Z$ be the poset on $\N$ whose only nonidentity arrows are
\[
 0\to1\leftarrow2\to3\leftarrow4\to\cdots.
\]
Equivalently there is an arrow $r\to s$ exactly when $r=s$, or
$r$ is even and $s$ is an adjacent odd integer.  No two nonidentity
arrows compose.  Let $\mathcal M$ be its free completion under
countable products.  Explicitly an object is a tuple
$x=(x_i)_{i\in I}$ with $I$ finite or countably infinite and each
$x_i\in\N$.  We choose the standard finite ordinals and $\N$ as
index sets; reindexing a countable disjoint union by one of them
makes the category small.  A morphism
\[
 x=(x_i)_{i\in I}\longrightarrow y=(y_j)_{j\in J}
\]
consists of a function $\phi:J\to I$ and, for every $j$, the
(necessarily unique when it exists) arrow $x_{\phi(j)}\to y_j$
in $Z$.  Composition composes the functions and those arrows.
The empty tuple is terminal.  A countable product is concatenation
of the tuples, with its coordinate projections; the displayed
description of morphisms verifies its universal property exactly.

Let $W$ consist of the morphisms whose coordinate function $\phi$
is bijective.  It contains identities and is closed under countable
products, because the disjoint union of bijections is a bijection.
It even has two out of three, though that is not needed here.

Call a tuple \emph{bounded} when all its integer labels are bounded
by one finite integer.  A forward morphism with nonempty source
bounded by $r$ has target bounded by $r+1$, since every arrow in
$Z$ changes its label by at most one.  The same statement holds
for a formal inverse of a member of $W$: its coordinate function
is bijective, and its coordinate arrows again change labels by
at most one.  The empty tuple cannot be followed by a forward
arrow to a nonempty tuple or by a backward $W$-arrow to a nonempty
tuple, so it cannot invalidate this argument.

Every morphism in a Gabriel--Zisman localization is represented
by a finite zigzag.  Starting at the singleton tuple $(0)$, every
object on any zigzag that can end at a nonempty object is therefore
bounded.  There is consequently no morphism
\[
 Q(0)\longrightarrow Q\bigl((n)_{n\in\N}\bigr).
\]
On the other hand, the singleton tuple $(n)$ is isomorphic to $(0)$
in the localization, by the first $n$ arrows of the displayed
zigzag, all of which belong to $W$.  Hence every set
$\Hom(Q(0),Q(n))$ is nonempty.  The actual countable product in
$\mathcal M$ of the singleton tuples $(n)$ is exactly the unbounded
tuple $(n)_{n\in\N}$.  Its localized cone cannot be a product:
the required map of hom sets has empty source and nonempty target.

The proof supplies the whole infinite category, every arrow and
composition, and a bound applying to every finite zigzag; it is
not an inference from bounded experiments.  The accompanying
program \texttt{code/G\_zigzag\_bound.py} verifies the sharp
finite-distance certificates for labels through $30$.  The
infinite argument is the uniform-bound argument just given.
\end{proof}

\begin{claim}{G.13}{Proved}{Infinite products with a calculus of fractions}
\label{G:fractions-products}
Let $I$ be any small set.  Suppose that $\mathcal M$ has
$I$-indexed products, that $W$ contains identities and is closed
under these products, and that $(\mathcal M,W)$ admits a left
or a right calculus of fractions.  Then $\mathcal M[W^{-1}]$
has $I$-indexed products and the localization functor preserves
the specified product cones.
\end{claim}
\begin{proof}
For a right calculus, every localized morphism is represented
by a roof $X\leftarrow X'\to Y$ whose left arrow is in $W$.
Two roofs represent the same morphism exactly when they have
a common refinement of roofs of the standard kind.  Therefore
the canonical functor
\[
 (\mathcal M^I)[(W^I)^{-1}]
       \longrightarrow (\mathcal M[W^{-1}])^I
\]
is full: choose a roof in each coordinate and assemble the
two coordinate families into one roof in $\mathcal M^I$.
It is faithful: assemble the coordinate common refinements into
one common refinement.  The calculus axioms themselves hold
coordinatewise, so the asserted criterion for equality is the
same on both sides.  Objects are identical.  A left calculus
uses roofs of the opposite orientation and gives the same
conclusion.  This is the uniform-length replacement for the
finite-padding argument of \ref{G:finite-products-general}.

Now the diagonal/product adjunction has relative left and right
functors.  Localize its unit and counit and use the displayed
equivalence.  The descended adjunction exhibits the product
and identifies its cone with the localization of the original
one.  Every map used here has a fixed length-two fraction
representative, so no unbounded family of zigzags occurs.
\end{proof}

\begin{claim}{G.14}{Proved}{Infinite products for model-category weak equivalences}
\label{G:model-products}
Let $\mathcal M$ be a model category, $W$ its weak equivalences,
and $I$ a small set.  If products indexed by $I$ preserve
families of weak equivalences, then its homotopy category has
those products and $\mathcal M\to\operatorname{Ho}(\mathcal M)$
preserves them.  Without this preservation hypothesis the
product of fibrant replacements computes the derived product,
and no assertion about the original product cone is intended.
\end{claim}
\begin{proof}
The product category $\mathcal M^I$, with all three classes
defined coordinatewise, is a model category.  Factorizations
and lifts are chosen coordinatewise, and its limits and
colimits exist coordinatewise.  Its cofibrant-fibrant objects
are the tuples of cofibrant-fibrant objects.  The usual
description of morphisms between such objects as homotopy
classes gives
\[
 \operatorname{Ho}(\mathcal M^I)
       \simeq\operatorname{Ho}(\mathcal M)^I.
\]
Here homotopies also assemble coordinatewise: choose the tuple
of cylinder or path objects in the product model category.
The homotopy relation is independent of those choices, so
coordinatewise equalities yield equality in the product.

The diagonal functor preserves weak equivalences.  By the
stated hypothesis so does the product functor.  Descend their
adjunction, as in \ref{G:fractions-products}, and use the
displayed equivalence to obtain the required products and
preservation.  Alternatively choose fibrant replacements
$X_i\to RX_i$.  Their product is a weak equivalence by
hypothesis, and $\prod_i RX_i$ is fibrant.  Maps from a
cofibrant replacement into this product are precisely tuples
of maps up to homotopy, which proves the same universal
property for the original product after localization.
\end{proof}

\paragraph{Manuscript conclusions.}
The assertion about finite products in XVI ms.~7--16 and
99--113 needs more than the literal weak axioms:
\ref{G:products-fail} gives its finite failure, and
\ref{G:no-left-adjoint} gives the failure of bare left Kan
extensions in XV ms.~134--137.  In the test class $\Wt$ the
nontriviality assertions of XV ms.~133--133 quinquies nevertheless
hold, as \ref{G:nontriviality-false} specifies.  For arbitrary
weak localizers they fail; the complete corrected statements
are \ref{G:ccprime} and \ref{G:retract-restoration}.
The direction-of-implication correction in R20 is confirmed.
The edition's doubt about infinite products at XVI Lemma~2.5
is justified by \ref{G:infinite-product-counterexample};
its model-category qualification is proved in
\ref{G:model-products}.  Forward gluing together with filtered
colimits restores relative Theorem~A by
\ref{G:gluing-recovery}, hence restores the familiar product
and Kan-extension assertions by \ref{G:restored-theory}.

